\chapter{Considerações Finais} \label{chap:conclusao}

Este capítulo encerra a dissertação. Retoma a questão de pesquisa e o percurso percorrido, avalia o cumprimento dos objetivos e a verificação das hipóteses, discute com franqueza as limitações do estudo, examina a possibilidade de generalização para outros ativos e aponta os desdobramentos futuros. A intenção é oferecer um balanço honesto, que reconheça tanto o que foi alcançado quanto o que permanece em aberto.

\section{Retomada do Problema e do Percurso}

A pesquisa partiu de uma questão precisa: de que forma a fusão do sentimento de notícias financeiras, extraído por um modelo \textit{Transformer} especializado em finanças, com a modelagem econométrica do risco pode contribuir para a previsão da direção e da volatilidade da PETR4. Para respondê-la, construiu-se uma arquitetura reprodutível que integra quatro componentes: a coleta de notícias com marcação temporal precisa, a organização do corpus por uma taxonomia temática supervisionada, a extração de sentimento pelo FinBERT-PT-BR e a modelagem do risco pelo GARCH, tudo reunido por fusão de dados e submetido a classificadores não lineares sob protocolo cronológico.

O contraste com a etapa de qualificação é o eixo desta versão. Naquela ocasião, a ausência de marcação temporal levou a uma imputação estocástica de datas que, como a banca observou com razão, desacoplava a notícia da reação do mercado e invalidava os resultados. A correção dessa falha, por meio da coleta direta via interfaces de programação que fornecem a data e a hora exatas, transformou um experimento de viabilidade em um estudo empírico legítimo, cujos resultados decorrem de um corpus real de duzentas e cinco mil notícias.

\section{Cumprimento dos Objetivos}

Os cinco objetivos específicos foram cumpridos. O corpus com marcação temporal foi coletado de cinco portais e alinhado pela regra \textit{Lead-Lag}. A taxonomia supervisionada de sete categorias e cento e cinquenta e dois termos substituiu, com vantagem de interpretabilidade, a modelagem de tópicos não supervisionada. O sentimento foi extraído por um modelo especializado no domínio financeiro, com o procedimento de obtenção do escore documentado em detalhe. A volatilidade foi modelada por um GARCH(1,1) cujo ajuste foi validado pela análise dos resíduos. Por fim, os classificadores foram treinados e avaliados sob divisão cronológica, com linhas de base e testes de significância.

\section{Verificação das Hipóteses}

A hipótese H1, do impacto direcional, foi confirmada de forma modesta. A fusão do sentimento com os preços elevou a acurácia de teste de cinquenta vírgula um para cinquenta e quatro vírgula cinco por cento, desempenho estatisticamente superior ao acaso pelo teste binomial, com valor de probabilidade de zero vírgula zero dois quatro. O ganho sobre o modelo de preços, contudo, tem significância marginal pelo teste de McNemar, o que situa a contribuição do sentimento em uma magnitude compatível com a quase-eficiência do mercado.

A hipótese H2, da antecipação de risco, foi sustentada de maneira robusta. A análise de causalidade de Granger mostrou que o sentimento antecede a volatilidade de forma forte e persistente em todas as defasagens testadas, evidência muito mais clara do que a obtida para a direção. Esse é, possivelmente, o achado mais relevante da dissertação, pois indica que o valor informacional do sentimento se manifesta sobretudo no canal de risco.

A hipótese H3, do viés de negatividade, foi corroborada. Além do forte predomínio de notícias negativas no corpus, a regressão quantílica forneceu evidência direta da assimetria: o efeito do sentimento sobre o retorno é grande e significativo na cauda inferior da distribuição, nos dias de pior desempenho, e desaparece nos quantis superiores. O sentimento das notícias importa, portanto, sobretudo nos momentos de queda, exatamente como prevê a assimetria comportamental dos investidores.

Convém um registro de método. Em contraste com a versão de qualificação, não se interpreta aqui a acurácia próxima de cinquenta por cento como prova de robustez. Aquela leitura decorria do ruído introduzido pela imputação de datas. Com os dados reais, a referência honesta é o ganho modesto, porém significativo, sobre as linhas de base.

\section{Contribuições da Pesquisa}

À luz dos resultados, é possível enunciar de forma precisa as contribuições deste trabalho, organizadas em quatro planos e sintetizadas na Tabela~\ref{tab:contribuicoes}.

No plano \textbf{teórico}, a principal contribuição é a quantificação da assimetria do efeito do sentimento. Por meio da regressão quantílica, demonstrou-se que o sentimento das notícias não exerce um efeito uniforme sobre o retorno, e sim um efeito concentrado na cauda inferior da distribuição, grande e estatisticamente significativo nos dias de pior desempenho e praticamente nulo nos dias de alta. Esse achado fornece evidência empírica direta do viés de negatividade para um ativo individual da B3, articulando o Teorema da Concordância e as Finanças Comportamentais, e estende ao nível de um ativo específico o que \citeonline{silva_o_2018} havia observado para o índice agregado.

No plano \textbf{metodológico}, a contribuição é dupla. A primeira é a consolidação de uma arquitetura unificada e reprodutível que integra, de forma inédita para a PETR4, um modelo de sentimento contextual de domínio, o FinBERT-PT-BR, o controle econométrico do risco pelo GARCH e a previsão por aprendizado de máquina, tudo sobre um corpus dotado de marcação temporal precisa, requisito que viabiliza a análise de antecedência entre notícia e preço. A segunda é a demonstração de que a combinação de regressões quantílicas com pesos variáveis melhora de forma substancial a previsão de volatilidade em relação ao modelo linear.

No plano \textbf{empírico}, a pesquisa entrega resultados reportados com rigor e transparência: o ganho do sentimento sobre o acaso na previsão de direção, estatisticamente significativo; a forte causalidade preditiva do sentimento sobre a volatilidade; e a delimitação honesta de que esse sinal, embora real, não se converte em melhoria de previsão pontual de volatilidade fora da amostra, conclusão que converge com a literatura e evita a superestimação.

No plano \textbf{prático}, o trabalho disponibiliza um \textit{pipeline} modular e transferível a outros ativos, uma distinção operacional entre eventos de mercado e de empresa por meio da taxonomia temática, e uma avaliação econômica que indica valor da vantagem preditiva no período analisado.

\begin{table}[htpb]
\centering
\caption{Síntese das contribuições da pesquisa e respectivas evidências.}
\label{tab:contribuicoes}
\resizebox{0.98\textwidth}{!}{%
\begin{tabular}{p{2.3cm} p{6.3cm} p{5.6cm}}
\hline
\textbf{Plano} & \textbf{Contribuição} & \textbf{Evidência} \\ \hline
Teórico & Quantificação da assimetria do efeito do sentimento (viés de negatividade) em um ativo individual & Efeito de $+542$ bps no quantil $0{,}05$ ($p=0{,}001$), nulo nos quantis altos \\
Metodológico & Arquitetura unificada com sentimento contextual, GARCH e ML, sob marcação temporal precisa & FinBERT-PT-BR; alinhamento Lead-Lag; divisão cronológica \\
Metodológico & Modelo quantílico ponderado para a previsão de volatilidade & $R^2$ fora da amostra de $+10{,}9\%$ sobre o HAR linear \\
Empírico & Ganho significativo do sentimento na previsão de direção & $54{,}5\%$ vs $50{,}1\%$; teste binomial $p=0{,}024$ \\
Empírico & O sentimento antecede a volatilidade & Causalidade de Granger $p \approx 0{,}000$ em todas as defasagens \\
Prático & Distinção mercado vs ativo e \textit{pipeline} transferível & Taxonomia de 7 categorias; generalização para a VALE3 \\ \hline
\end{tabular}%
}
\vspace{0.2cm}
{\small \\ Fonte: Elaborado pelo autor.}
\end{table}

Em síntese, a contribuição central da dissertação pode ser assim enunciada: trata-se de uma arquitetura reprodutível e honesta que, pela primeira vez, aplica um modelo de sentimento contextual de domínio a um ativo individual da B3 com marcação temporal precisa, e cujo achado mais relevante é o de que o efeito do sentimento sobre o retorno não é uniforme, mas concentrado nos momentos de queda, resultado que a análise na média e a previsão de direção, isoladamente, não revelam.

\section{Limitações da Pesquisa}

A credibilidade de um estudo depende do reconhecimento de suas limitações, e esta dissertação assume as seguintes. A magnitude do sinal direcional é modesta, da ordem de cinquenta e quatro por cento, de modo que as previsões têm caráter acadêmico e não constituem recomendação de investimento. O sentimento é extraído do título e do resumo das notícias, e não do texto integral, e a volatilidade realizada é aproximada pelo valor absoluto do retorno, uma medida reconhecidamente ruidosa. O FinBERT-PT-BR foi aplicado sem reajuste ao corpus específico da pesquisa, e a validação humana por coeficiente \textit{kappa}, embora preparada, ainda será consolidada. O estudo de refinamento revelou sensibilidade a mudanças de regime entre períodos, o que limita a generalização temporal dos modelos mais complexos. Por fim, a análise restringe-se a um único ativo e a um único idioma.

Além dessas limitações, cabe apontar ameaças à validade, organizadas segundo as categorias usuais. A validade externa, ou capacidade de generalização, é limitada pela escolha de um único ativo e de um único período, o que será discutido na próxima seção; resultados obtidos para a PETR4, ativo de características idiossincráticas, não se transferem automaticamente a outros papéis. A validade de construto depende da adequação da polaridade das manchetes como aproximação do sentimento de mercado; embora amparada pela literatura, essa aproximação ignora o corpo da notícia e a sua repercussão, e o próprio FinBERT-PT-BR, treinado em um corpus financeiro genérico, pode não capturar nuances específicas da cobertura sobre a Petrobras. A validade interna, isto é, a garantia de que a relação medida não decorre de artefatos, é protegida pela divisão estritamente cronológica e pela defasagem de todos os atributos, que evitam o vazamento de informação futura, mas permanece sujeita ao risco de variáveis omitidas, como fatores macroeconômicos não modelados que afetem simultaneamente o sentimento e os preços. A validade de conclusão estatística, por fim, é resguardada pelo uso de testes formais de significância e pelo relato de intervalos de confiança, ainda que o tamanho da amostra de teste imponha limites ao poder estatístico, especialmente nas análises por subperíodo e por regime. O reconhecimento explícito dessas ameaças não enfraquece o trabalho; ao contrário, delimita com honestidade o alcance de suas conclusões, postura que a banca de qualificação expressamente valorizou.

\section{Generalização para Outros Ativos}

Uma das questões levantadas pela banca diz respeito à possibilidade de estender o estudo a outros ativos, como a VALE3. A arquitetura foi desenhada de forma modular e é, em boa medida, transferível. A migração exigiria três ajustes. O primeiro é a substituição da série financeira-alvo pelas cotações do novo ativo e a reestimação do GARCH. O segundo é a adaptação da taxonomia temática ao setor correspondente, de modo que, no caso da VALE3, o vocabulário de petróleo e refino daria lugar ao de minério de ferro, à China como principal destino de exportação e a questões ambientais e de barragens. O terceiro é o ajuste dos termos de busca da coleta. A camada de processamento de linguagem, baseada no FinBERT-PT-BR, e a arquitetura de fusão de dados permaneceriam inalteradas. Essa modularidade é, em si, uma contribuição prática do trabalho, pois indica um caminho de replicação de baixo custo.

Vale antecipar, contudo, que a transposição não seria mecânica, e que a sua análise traria interesse próprio. A VALE3, embora também sensível a uma \textit{commodity} e a fatores políticos, difere da PETR4 na natureza da distinção entre mercado e ativo: para uma exportadora de minério de ferro, um aquecimento da demanda chinesa é simultaneamente uma boa notícia para o preço da \textit{commodity} e para a empresa, sem a inversão de sinal que caracteriza os choques de oferta no petróleo. Já eventos como o rompimento de barragens introduzem uma dimensão de risco operacional e reputacional ausente no caso da Petrobras. Comparar as duas aplicações permitiria, assim, testar se a assimetria do efeito do sentimento, identificada para a PETR4, é uma característica geral do mercado brasileiro ou um traço específico de ativos de determinado perfil. Essa investigação comparativa, mais do que uma simples replicação, configura uma agenda de pesquisa promissora para o doutorado.

\section{Trabalhos Futuros}

O estudo delimita com clareza onde não está o ganho, que não reside em mais atributos nem em mais ajuste do classificador, e aponta direções mais promissoras. A primeira é o ajuste fino do FinBERT-PT-BR ao corpus da pesquisa, a partir do conjunto-ouro rotulado, com vistas a elevar a qualidade do sinal de sentimento na origem. A segunda é a previsão da magnitude da volatilidade, e não apenas da direção. A dissertação já implementou, com esse fim, a regressão quantílica da volatilidade com pesos variáveis atribuídos aos quantis, que melhorou de forma substancial a previsão em relação ao modelo linear, embora tenha revelado que o ganho provém da própria metodologia, e não do sentimento. O caminho futuro é, portanto, buscar extrair o sinal de volatilidade do sentimento por meio de proxies de volatilidade de alta frequência e de modelos sensíveis a regime, capazes de captar a contribuição que a frequência diária e a abordagem aqui adotada não revelaram. A terceira é a exploração de janelas intradiárias, que aproveite a marcação temporal precisa para medir a reação do mercado em horizontes curtos. A quarta é a adoção de modelos de volatilidade sensíveis a regime e de especificações assimétricas, como o EGARCH, que capturem o efeito de alavancagem. A quinta é o aprofundamento da avaliação da significância econômica. A simulação de negociação preliminar reportada no Capítulo~\ref{chap:resultados} indicou que a vantagem preditiva supera a estratégia passiva no período de teste, líquida de custos, mas a sua extensão a múltiplos regimes e a diferentes estruturas de custo permanece como tarefa futura. A sexta é a generalização multiativo, com a replicação do \textit{pipeline} para a VALE3 e outros papéis de alta liquidez.

\section{Discussão Geral dos Achados}

Reunidos, os resultados desta dissertação compõem um quadro coerente sobre o papel do sentimento das notícias na precificação da PETR4, quadro que convém discutir de forma integrada. O ponto de partida é que o sinal do sentimento existe, mas é modesto e sutil, exatamente como prevê a quase-eficiência do mercado. A previsão de direção, ainda que estatisticamente superior ao acaso, situa-se pouco acima do patamar de cinquenta por cento, e o ganho marginal sobre o modelo de preços é pequeno. Tomado isoladamente, esse resultado poderia sugerir que o sentimento tem pouca relevância prática.

A análise mais fina, contudo, revela que essa leitura seria precipitada, e é aqui que reside a contribuição mais original do trabalho. Ao examinar o efeito do sentimento não na média, mas ao longo da distribuição do retorno, descobre-se que ele é fortemente assimétrico: concentra-se nos dias de queda, onde se torna grande e significativo, e desaparece nos dias de alta. Esse padrão concilia, em um único achado empírico, três fios teóricos desenvolvidos no referencial. Conecta-se ao viés de negatividade das Finanças Comportamentais, segundo o qual as más notícias pesam mais; à teoria dos limites à arbitragem, que explica por que o sentimento pode mover os preços sem ser corrigido de imediato; e à evidência de que o efeito da mídia se intensifica nos momentos de tensão. O sentimento, portanto, importa, mas importa onde mais importa, isto é, nas quedas, que são precisamente os eventos de maior relevância para a gestão de risco.

No tocante à volatilidade, o estudo traça uma fronteira nítida entre associação e previsão. Há associação forte e robusta: o sentimento antecede a volatilidade de maneira persistente, como atesta a causalidade de Granger, e a sua intensidade é um preditor significativo dentro da amostra. Não há, porém, ganho de previsão pontual fora da amostra, mesmo quando se emprega a modelagem quantílica ponderada que se mostrou superior para a volatilidade. Essa distinção, longe de ser um fracasso, é um esclarecimento de valor, pois delimita com precisão o alcance do sinal e converge com o que a literatura nacional havia encontrado. Por fim, a sensibilidade dos resultados ao regime de mercado, evidenciada tanto pela análise por subperíodo quanto pelo condicionamento à incerteza, confirma que o efeito do sentimento é dependente do estado do mercado, e não uma constante universal. O conjunto desses achados, reportado sem exagero, oferece um retrato realista e teoricamente ancorado do fenômeno estudado.

\section{Considerações de Encerramento}

A dissertação demonstrou, com dados reais e rigor estatístico, que o sentimento de notícias financeiras contribui de forma modesta, porém significativa, para a previsão da direção da PETR4, e de forma mais nítida para a sua volatilidade. Ao corrigir a falha de marcação temporal, remover componentes injustificados, adotar um modelo de sentimento especializado e reportar os resultados sem superdimensionamento, o trabalho consolidou uma base metodologicamente sólida e reprodutível. Mais do que um número de acurácia, o legado do estudo é uma arquitetura honesta e transferível para investigar o papel informacional da mídia na precificação de ativos, alinhada às recomendações da banca de qualificação e preparada para a defesa definitiva e para os seus desdobramentos no doutorado.
